\documentclass{article}
\usepackage[T1]{fontenc}
\usepackage{lmodern}
\usepackage{graphicx}
\usepackage{amsmath,amssymb}
\usepackage[a4paper,margin=2cm]{geometry}
\usepackage[absolute,overlay]{textpos}
\setlength{\TPHorizModule}{1mm}
\setlength{\TPVertModule}{1mm}
\usepackage[indonesian]{babel}
\usepackage{hyperref}
\setlength{\emergencystretch}{2em}
% unit: Halaman 61

\begin{document}

% === Halaman 61 (python/native) ===

• Cipher berulang dinyatakan sebagai




𝐶𝑖 = 𝑓(𝐶𝑖−1, 𝐾𝑖) 



    yang dalam hal ini,


i = 1, 2, …, r    (r adalah jumlah putaran). 


Ki = kunci putaran (round key) pada putaran ke-i


f = fungsi transformasi (round function), di dalamnya terdapat  operasi

                   substitusi, permutasi, dan/atau ekspansi, kompresi). 


• Plainteks  dinyatakan dengan C0 dan cipherteks dinyatakan dengan Cr.

61

\end{document}
